\section{Related Work}
\label{sec:related}

\subsection{Message Passing and Expressive Power}

Modern graph learning descends from early recursive models on graph
domains~\cite{gori2005new} and the spectral convolution of
Kipf and Welling~\cite{kipf2017gcn}. GraphSAGE introduced inductive
neighborhood aggregation~\cite{hamilton2017inductive}, graph attention networks
weighted neighbors by learned attention~\cite{velivckovic2017graph}, and the
message-passing neural network framework unified these designs for molecular
property prediction~\cite{gilmer2017neural}. The Weisfeiler--Lehman connection
established by GIN characterizes the discriminative power of this family and
shows that injective multiset aggregation is central to its
expressiveness~\cite{xu2019gin}. Surveys document the resulting design
space~\cite{wu2020comprehensive}, and more recent attention-based
variants extend it further~\cite{sun2023attention}.

\subsection{Deep GNNs and Residual Connectivity}

Depth is not free in message passing. Over-smoothing makes node
representations converge as propagation depth grows~\cite{li2018deeper}, and
this limiting behavior has been analyzed
explicitly~\cite{oono2020simple,chenSimpleDeepGraph2020}. Over-squashing is the
complementary failure mode, in which long-range information is compressed into
fixed-size messages~\cite{alon2020bottleneck,topping2021understanding}. Several
mechanisms target these effects: DropEdge randomly removes edges during
training~\cite{rong2019dropedge}, Jumping Knowledge aggregates representations
across layers~\cite{xu2018representation}, and residual or dense skip paths
carry earlier states forward~\cite{he2016deep,li2019deepgcns,li2021training}.
Residual gated graph convolutional networks learn how strongly historical
states should be reinjected~\cite{bresson2017residual}, and DenseGNN stacks
dense connections for scalable deep
prediction~\cite{du2024densegnn}. These approaches reuse depth-wise state
within a single stream. Our work keeps the same reuse principle but asks a
different question: given several parallel branches, which local residual
neighborhoods on the branch-by-layer grid should be connected.

\subsection{Multi-Branch and Higher-Order Architectures}

Higher-order neighborhood mixing shows that propagating over several hops with
sparsified mixing improves over single-hop
aggregation~\cite{abu2019mixhop}. ARMA filters supply a related frequency-domain
view of multi-hop propagation~\cite{bianchi2020graph}, and graph transformers
generalize attention beyond local
neighborhoods~\cite{dwivedi2020generalization}. Hierarchical pooling adds a
second axis of structure by coarsening graphs as they are
processed~\cite{ying2018hierarchical,lee2019self,cangea2018towards}, and recent
surveys catalogue the resulting
designs~\cite{li2024graph,sun2023attention}. These lines of work increase the
expressive or aggregational capacity of the propagation layer. In contrast, we
hold the propagation operator fixed and vary only how hidden states are reused
across a branch-by-layer grid, so that the residual topology itself becomes the
identifiable factor.

\subsection{Controlled Evaluation for Graph Classification}

Because graph benchmarks are sensitive to splitting, feature choice, and
operator tuning, careful comparative protocols matter. Errica \emph{et al.}
argue for matched evaluation conditions when comparing graph
classifiers~\cite{errica2020fair}, and TUDataset provides the standardized
collection on which much of this evaluation
rests~\cite{morris2020tudataset}. We follow the same discipline: the readout,
optimizer schedule, fold partition, and backbone options are matched across the
compared residual families, and only the residual neighborhood and
residual-control mechanism change.
